\documentclass[11pt,a4paper]{article}
\usepackage[margin=2.25cm]{geometry}
\usepackage{amsmath,amssymb,mathtools,booktabs,microtype}
\usepackage{graphicx}
\usepackage[hidelinks]{hyperref}
\usepackage{enumitem}
\setlength{\parindent}{0pt}
\setlength{\parskip}{5pt}
\makeatletter
\renewcommand\@seccntformat[1]{\csname the#1\endcsname\quad}
\makeatother
\newcommand{\Tr}{\operatorname{Tr}}
\newcommand{\norm}[1]{\left\lVert #1\right\rVert}
\newcommand{\figurefile}[2]{%
  \IfFileExists{#1}{\includegraphics[width=#2\linewidth]{#1}}%
  {\fbox{\parbox[c][0.22\textheight][c]{#2\linewidth}{\centering
  Run the notebook export cell to create this figure.}}}}
\title{Kernel Polynomial Hartree--Fock with Structured Hadamard Probes}
\author{Method note for the large-scale quasiperiodic calculations}
\date{August 2026}
\begin{document}
\maketitle

\begin{abstract}
This note describes the kernel polynomial method (KPM) implementation used to
obtain zero-temperature density matrices and local Hartree--Fock fields on
large open quasiperiodic lattices. The method never stores the density matrix:
it applies an approximation to the occupied projector to blocks of vectors and
estimates only the diagonal and nearest-neighbour elements needed by the
Hartree and Fock terms. We explain the polynomial, chemical-potential, and
self-consistent parts of the calculation; the limitations of ordinary random
trace estimators; and why randomized, nested Hadamard probes are the practical
choice for the present accuracy target.
\end{abstract}

\section{\hspace{0.5em}Mean-field problem}
We consider a real, spinless nearest-neighbour Hamiltonian on an open square
lattice. The hopping amplitudes may be quasiperiodically modulated. At a
Hartree--Fock fixed point the effective one-particle Hamiltonian is
\begin{equation}
 H_{\rm eff}=H_0+V_{\rm H}+V_{\rm F},\qquad
 \rho_{ij}=\langle c_j^\dagger c_i\rangle .
\end{equation}
For nearest-neighbour repulsion $U$, the local fields used in the calculation
are
\begin{align}
 (V_{\rm H})_{ii} &= U\sum_{j\in {\rm nn}(i)}\rho_{jj}, \\
 (V_{\rm F})_{ij} &= -U\,\mathrm{Re}\,\rho_{ij},
 \qquad i,j\text{ nearest neighbours}.
\end{align}
Only $n_i=\rho_{ii}$ and the two forward bond fields are consequently needed;
a dense $N\times N$ density matrix is unnecessary. At zero temperature and
fixed particle number $N_e$, the target is
\begin{equation}
 P=\Theta(\mu-H_{\rm eff}),\qquad \Tr P=N_e .
 \label{eq:projector}
\end{equation}
KPM avoids constructing $P$ by applying an approximation to it to probe
vectors. The chemical potential is adjusted at every SCF iteration so that the
approximated trace has the requested filling.

\section{\hspace{0.5em}Chebyshev approximation}
An enclosing spectral interval $[E_-,E_+]$ is required for each effective
Hamiltonian. With
\begin{equation}
 \widetilde H=\frac{H_{\rm eff}-bI}{a},\qquad
 a=\frac{E_+-E_-}{2},\qquad b=\frac{E_++E_-}{2},
\end{equation}
the spectrum lies in $[-1,1]$. The chemical potential is similarly scaled as
$\widetilde\mu=(\mu-b)/a$. A degree-$M$ KPM projector is
\begin{equation}
 P_M(\mu)=\sum_{m=0}^{M}g_m c_m(\widetilde\mu) T_m(\widetilde H),
 \label{eq:kpm}
\end{equation}
where $T_m$ is a Chebyshev polynomial, $c_m$ is the occupied-step coefficient,
and $g_m$ is a damping-kernel coefficient. Kernel damping is essential: a
bare truncated representation of a discontinuous Fermi step has Gibbs
oscillations, giving fractional occupations and oscillatory local fields.
Finite $M$ is therefore a controlled spectral broadening, not an exact sharp
projector.

For a vector $z$, no matrix polynomial is formed. Instead,
\begin{align}
 q_0&=z, &q_1&=\widetilde H z,\\
 q_{m+1}&=2\widetilde Hq_m-q_{m-1},& P_Mz&=\sum_m g_mc_mq_m .
 \label{eq:recurrence}
\end{align}
Only three vector blocks are live during the recurrence. The nearest-neighbour
Hamiltonian has a bounded local stencil, so applying it to a block of $B$
probes costs $\mathcal O(NB)$. The dominant cost is $\mathcal O(NMR)$ for $R$
probes. Blocking makes this compatible with GPU memory, which is why the
method can treat millions of sites without a density matrix in memory.

\section{\hspace{0.5em}Trace correction and local estimators}
Let $Z=[z_1,\ldots,z_R]$ and $Y=P_MZ$. The estimators used for the trace,
density, and bond order are
\begin{align}
 \widehat{\Tr P_M}&=\frac1R\sum_r z_r^\mathsf T y_r,\\
 \widehat n_i&=\frac1R\sum_r z_{ir}y_{ir},\\
 \widehat\rho_{ij}&=\frac1R\sum_r z_{ir}y_{jr}.
 \label{eq:estimators}
\end{align}
The last estimator is required only on horizontal and vertical bonds. The
Hartree and Fock fields are then rebuilt directly from these local arrays. The
trace equation
\begin{equation}
 \widehat N(\mu)=\widehat{\Tr P_M(\mu)}=N_e
\end{equation}
is solved for $\mu$ during every SCF step. This is not optional at high
precision: finite polynomial order, open boundaries, and updated fields all
shift the apparent Fermi level.

\section{\hspace{0.5em}What goes wrong with ordinary stochastic probes}
For independent Rademacher vectors, $z_i=\pm1$, the estimators in
Eq.~\eqref{eq:estimators} are unbiased. But their local variance is
\begin{equation}
 \mathrm{Var}(\widehat n_i)=\frac1R\sum_{j\ne i}|(P_M)_{ij}|^2 .
 \label{eq:variance}
\end{equation}
Thus local noise is controlled by off-diagonal projector weight, not by the
trace alone. In a low-gap or strongly quasiperiodic problem that weight can be
long ranged and spatially nonuniform.

Three limitations make independent random probes a poor production default for
self-consistent local fields. First, noise in $n_i$ and $\rho_{ij}$ enters the
Hartree potential and Fock hopping, so it becomes noise in the SCF map itself.
Apparent stagnation may then be sampling noise rather than a physical failure
of fixed-point iteration. Second, an accurate total trace does not certify an
accurate localized density peak or bond: scalar errors can cancel. Third,
independent $R$ and $2R$ samples do not give a clean convergence comparison,
because the probe subspace and random realization both change.

The usual $R^{-1/2}$ behaviour is costly. A tenfold reduction of the standard
error requires roughly one hundred times more independent probes, and that
cost is incurred at every SCF step. Random probes remain useful for rough
estimates, but they do not give the reproducible local error control required
here.

\section{\hspace{0.5em}Randomized Hadamard probes}
For $N=2^L$, let $H_N$ be a Walsh--Hadamard matrix. It satisfies
\begin{equation}
 H_NH_N^\mathsf T=NI .
\end{equation}
The calculation selects columns of a randomly signed and randomly permuted
Hadamard basis. Each probe has only $\pm1$ entries, is generated efficiently
from bit parity, and is exactly orthogonal to every other selected probe. At
full coverage,
\begin{equation}
 \frac1N H_NH_N^\mathsf T=I.
 \label{eq:hadamardexact}
\end{equation}
Consequently, for $R=N$, all estimators in Eq.~\eqref{eq:estimators} are exact
for \emph{any} matrix $P_M$: the trace, every density, and every selected bond
have no probe error. A finite Rademacher sample has no analogous endpoint.

At $R<N$, Hadamard probing is not magically exact and does not remove finite
polynomial error. Its operational advantages are nevertheless decisive:
\begin{itemize}[leftmargin=1.7em]
\item orthogonality avoids redundant probe directions;
\item nested prefixes $R\subset2R\subset4R\cdots$ give a directly interpretable
same-realization probe ladder;
\item powers of two align with the bit-quantized lattice indexing;
\item a fixed scrambling produces a deterministic, reproducible SCF map; and
\item a second scrambling gives an independent probe check.
\end{itemize}
Hadamard probing is not a mathematical requirement of KPM. It is necessary in
the practical sense relevant here: it provides reproducible local observables,
an exact full-coverage limit, and a meaningful $R\to2R$ uncertainty test at a
probe count far below $N$.

\section{\hspace{0.5em}Separating the three numerical errors}
There are three distinct approximations:
\begin{enumerate}[leftmargin=1.7em]
\item SCF error: input and output fields have not reached a fixed point;
\item polynomial error: degree $M$ still broadens the Fermi projector; and
\item probe error: finite $R$ estimates local observables imperfectly.
\end{enumerate}
They must be measured separately. On a dense-reference system, calibrate both
$M$ and $R$ against exact local density and bonds. On a large system, save a
fixed effective field and compare nested $R$ and $2R$ at fixed $M$; then
compare $M$ with higher $M$ at fixed nested $R$. Only after those fixed-field
tests should one choose production parameters for full SCF. Finally, rebuild
the converged fields and perform an independent higher-$M$ audit with a second
Hadamard scrambling. Report RMS and maximum local density and bond differences,
not only energy and particle number.

The KPM global idempotency diagnostic is
\begin{equation}
 d_{\rm tr}=\frac{|\widehat{\Tr(P_M^2)}-\widehat{\Tr(P_M)}|}
 {|\widehat{\Tr(P_M)}|},
\end{equation}
where $\Tr(P_M^2)$ is estimated from $\norm{P_Mz_r}^2$. It diagnoses global
fractional occupation. It is useful but is not the MPO/SP2 relative Frobenius
residual $\norm{P^2-P}_F/\norm P_F$, and it cannot bound the largest local
error. The independent audit is therefore essential.

\section{\hspace{0.5em}Computed strong-quasiperiodic example}
The $256\times256$ strong quasiperiodic calculation used $M=40000$, $R=2048$
inside SCF and an independent $M=60000$, $R=2048$ audit. It reached the
relative $10^{-3}$ SCF threshold in nine iterations. Production and audit
agree in checkerboard order to $4.7\times10^{-8}$, in total energy to about
$2\times10^{-3}$, and in density RMS to $1.86\times10^{-5}$. The maximum
sitewise density difference is $4.49\times10^{-3}$, showing why global
observables alone are not a sufficient accuracy certificate in a localized
regime. The trace idempotency defect was about $1.5\times10^{-5}$ for both
calculations.

The accompanying notebook exports the figures below from the saved density
maps. The Fourier plot uses a Tukey window to suppress open-boundary
background before the FFT; it is therefore a diagnostic of bulk modulation,
not a periodic-boundary structure factor.

\begin{figure}[t]
\centering
\figurefile{figures/strong_m40000/density_full.pdf}{.92}
\caption{Full production density relative to half filling.}
\label{fig:density}
\end{figure}
\begin{figure}[t]
\centering
\figurefile{figures/strong_m40000/density_audit_difference.pdf}{.98}
\caption{Audit density and the local production--audit difference near the
largest discrepancy.}
\label{fig:audit}
\end{figure}
\begin{figure}[t]
\centering
\figurefile{figures/strong_m40000/density_fft_tukey.pdf}{.92}
\caption{Tukey-windowed density Fourier spectrum, with expected
quasiperiodic sidebands marked by crosses.}
\label{fig:fft}
\end{figure}

\section{\hspace{0.5em}Scope and limitations}
KPM with structured probes changes the dominant task from density-matrix
storage to repeated sparse local Hamiltonian applications. It can therefore
reach lattices far beyond dense diagonalization. It does not remove the need
for calibration: small gaps demand high $M$; rare localized features can make
maximum errors much larger than RMS errors; and a new hopping pattern,
interaction range, or boundary condition requires its own $M$/$R$ evidence.
The method is trustworthy only when the SCF, polynomial, and probe errors are
reported separately.
\end{document}
